\documentclass[11pt,a4paper]{article}

% ── Packages ──────────────────────────────────────────────
\usepackage[margin=2cm]{geometry}
\usepackage{graphicx}
\usepackage{listings}
\usepackage{xcolor}
\usepackage{hyperref}
\usepackage{caption}
\usepackage{subcaption}
\usepackage{float}
\usepackage{enumitem}
\setlist{nosep}

% ── Code listing style ────────────────────────────────────
\lstset{
  language=VHDL,
  basicstyle=\ttfamily\scriptsize,
  keywordstyle=\color{blue}\bfseries,
  commentstyle=\color{gray},
  numbers=left,
  numberstyle=\tiny\color{gray},
  frame=single,
  breaklines=true,
  captionpos=b,
}

% ── Title ─────────────────────────────────────────────────
\title{\textbf{Assignment 3: Interfaces and Memories}\\
       \large DIG 2026}
\author{
  Group X \\[4pt]
  Name 1 — Student ID \\
  Name 2 — Student ID \\
  Name 3 — Student ID
}
\date{\today}

% ══════════════════════════════════════════════════════════
\begin{document}
\maketitle

% ── 1. Component Description ─────────────────────────────
\section{Component Description}

% TODO: Describe each VHDL module / component used:
%  - UART RX (receiver)
%  - UART TX (transmitter)
%  - RAM (data buffer between RX and TX on FPGA1)
%  - Top-level entity for FPGA1 and FPGA2
%  - Baud-rate generator / clock divider
%  - (Bonus) SPI master/slave modules

\begin{table}[H]
\centering
\caption{System components overview}
\begin{tabular}{|l|l|p{6cm}|}
\hline
\textbf{Component} & \textbf{Location} & \textbf{Role} \\ \hline
UART RX        & FPGA 1 \& 2 & Receives serial data from PC / other FPGA \\ \hline
UART TX        & FPGA 1 \& 2 & Transmits serial data to other FPGA / PC  \\ \hline
RAM            & FPGA 1      & Stores received message before forwarding  \\ \hline
Baud Generator & FPGA 1 \& 2 & Generates timing for UART at target baud rate \\ \hline
% Add more rows as needed
\end{tabular}
\end{table}


% ── 2. Circuit Operation ─────────────────────────────────
\section{Circuit Operation}

% TODO: Explain the main block design and interactions.
%       Include a block diagram (place in figures/).

% \begin{figure}[H]
%   \centering
%   \includegraphics[width=0.9\textwidth]{figures/block_diagram.png}
%   \caption{Top-level block diagram of the two-FPGA system.}
%   \label{fig:block}
% \end{figure}


% ── 3. Data Flow Analysis ────────────────────────────────
\section{Data Flow Analysis}

The complete data path is:

\begin{center}
\texttt{PC1 $\rightarrow$ UART1 $\rightarrow$ FPGA\,1\,(RX) $\rightarrow$ RAM $\rightarrow$ FPGA\,1\,(TX) $\rightarrow$ UART1--2 $\rightarrow$ FPGA\,2\,(RX) $\rightarrow$ FPGA\,2\,(TX) $\rightarrow$ UART2 $\rightarrow$ PC2}
\end{center}

% TODO: Explain each stage in 3–5 sentences.


% ── 4. Simulation Results ────────────────────────────────
\section{Simulation Results}

\subsection{Functional Simulation}
% TODO: Insert waveform screenshot and explain.
% \begin{figure}[H]
%   \centering
%   \includegraphics[width=0.9\textwidth]{figures/func_sim.png}
%   \caption{Functional simulation waveform.}
%   \label{fig:func_sim}
% \end{figure}

\subsection{Post-Synthesis Simulation}
% TODO: Insert waveform screenshot and compare with functional.
% \begin{figure}[H]
%   \centering
%   \includegraphics[width=0.9\textwidth]{figures/post_synth_sim.png}
%   \caption{Post-synthesis simulation waveform.}
%   \label{fig:post_synth}
% \end{figure}

\subsection{Implementation Timing Simulation}
% TODO: Insert waveform screenshot and compare with previous two.
% \begin{figure}[H]
%   \centering
%   \includegraphics[width=0.9\textwidth]{figures/timing_sim.png}
%   \caption{Implementation timing simulation waveform.}
%   \label{fig:timing_sim}
% \end{figure}

\subsection{Comparison}
% TODO: 3–5 sentences comparing the three simulation types.
%       - Functional: ideal, no delays
%       - Post-synthesis: includes gate-level delays
%       - Timing: includes routing delays — closest to real hardware


% ── Bonus: SPI ────────────────────────────────────────────
% \section*{Bonus: SPI Communication}
% TODO: If attempted, describe the SPI link replacing UART1–2.


% ── Video Demo ────────────────────────────────────────────
\section*{Video Demonstration}
% TODO: Insert link
Video link: \url{https://youtube.com/your-video-link-here}

\end{document}
